% Generated draft: requires mathematical and source review.
\documentclass[11pt]{article}
\usepackage[margin=25mm]{geometry}
\usepackage{fontspec}
\setmainfont{Latin Modern Roman}
\usepackage{amsmath,amssymb,mathtools,unicode-math}
\setmathfont{Latin Modern Math}
\usepackage{enumitem,needspace,xurl,hyperref}
\setlength{\parindent}{0pt}
\setlength{\parskip}{5pt}
\setlength{\emergencystretch}{3em}
\newcommand{\N}{\mathbb{N}}
\newcommand{\Z}{\mathbb{Z}}
\newcommand{\Q}{\mathbb{Q}}
\newcommand{\R}{\mathbb{R}}
\newcommand{\C}{\mathbb{C}}
\newcommand{\F}{\mathbb{F}}
\newcommand{\A}{\mathbb{A}}
\newcommand{\PP}{\mathbb{P}}
\newcommand{\E}{\mathbb{E}}
\newcommand{\eps}{\varepsilon}
\newcommand{\dd}{\,\mathrm{d}}
\begin{document}
\section*{UnsolvedMath ID: 30005660}
\textbf{Source problem code:} catalogue entry
\subsection{Definitions and mathematical statement}
Let \(\mathcal T\) be a finite tensor category, in the convention intended by the supplied source. The record does not specify the ground field, its characteristic, or the precise conventions included in this term. Its stable category \(\underline{\mathcal T}\) has the same objects as \(\mathcal T\), with
\[
\operatorname{Hom}_{\underline{\mathcal T}}(X,Y)
=
\operatorname{Hom}_{\mathcal T}(X,Y)\big/
\{f:X\to Y\mid f\text{ factors through a projective object}\}.
\]
Write \(\mathbf 1\) for the tensor unit. The usual cohomology ring is
\[
H^\bullet(\mathcal T)=\bigoplus_{n\geq 0}\operatorname{Ext}_{\mathcal T}^{n}(\mathbf 1,\mathbf 1),
\]
with Yoneda multiplication. The supplied statement denotes the categorical center of this cohomology ring by \(C_{\mathcal T}^{\bullet}\), but supplies no definition of this particular center. Its precise construction must therefore be recovered from the cited source; the terminology alone does not identify it with the graded center of a ring or with a Drinfeld center.

For a nonnegatively graded commutative ring \(C^\bullet\), \(\operatorname{Proj}C^\bullet\) denotes the homogeneous prime ideals not containing the irrelevant ideal \(C_+=\bigoplus_{n>0}C^n\), with the Zariski topology. The analogous convention appropriate to the intended graded ring \(C_{\mathcal T}^{\bullet}\) is required here. The notation \(\operatorname{Spc}(\underline{\mathcal T})\) denotes the source's Balmer spectrum of the stable category. Because the statement does not assume a symmetric or braided tensor structure, the intended notion of prime thick tensor ideal and its spectrum must be specified from the source; the ordinary symmetric tensor-triangular convention cannot silently be imposed.

The supplied record asserts the existence of a continuous central-support comparison map
\[
\rho:\operatorname{Spc}(\underline{\mathcal T})\longrightarrow\operatorname{Proj}C_{\mathcal T}^{\bullet}.
\]
Its construction is not included in the record.

The source conjecture is that, for every finite tensor category \(\mathcal T\), this map \(\rho\) is a homeomorphism. The universal quantifier is retained without adding finite-generation, symmetry, braiding, or other hypotheses. This is a source-relative formulation: the missing conventions, definition of \(C_{\mathcal T}^{\bullet}\), and construction and well-definedness of \(\rho\) prevent a fully self-contained precise formulation from the supplied record alone.
\subsection{Short English statement}
The supplied conjecture asks whether the central-support comparison map identifies the spectrum of the stable category of every finite tensor category with the projective spectrum of its categorical cohomological center, including their topologies. The definitions of that center, the comparison map, and the spectrum convention for tensor categories without assumed symmetry require clarification. The record labels the problem open; this status is supplied curation, not an independently verified conclusion.
\subsection{Sources}
\begin{itemize}
\item Supplied local UnsolvedMath registry record, permanent ID \(30005660\), title \emph{Spectrum Comparison for Finite Tensor Categories}, exact source problem code \texttt{catalogue entry}, in the Oberwolfach Reports Open Problems collection. No individual proposer is identified. The proposed year is recorded as 2023.
\item Supplied report citation: \emph{Tensor-Triangular Geometry and Interactions (2024)}, Oberwolfach Reports, DOI \texttt{10.4171/owr/2023/40}, pp.\ 2303--2358. \url{https://doi.org/10.4171/owr/2023/40}. The record attributes the conjecture to the report's workshop discussion; it gives no specific page or speaker for this problem. The report year is reproduced as supplied.
\item Nakano--Vashaw--Yakimov, \emph{On the spectrum and support theory}, as listed in the supplied literature sources. \url{https://arxiv.org/abs/2112.11170}.
\item Vashaw, \emph{Balmer spectra and Drinfeld centers}, as listed in the supplied literature sources. \url{https://doi.org/10.2140/ant.2024.18.1081}.
\item Literature assessment dated exactly \texttt{2026-08-22}: the supplied curation labels the problem open and states that the comparison is a homeomorphism under stronger hypotheses and in many examples, while the assertion for every finite tensor category remains conjectural. These are supplied dated curation claims, not independently verified here. The assessment refers to a homogeneous spectrum, whereas the mathematical statement specifically uses \(\operatorname{Proj}\); their intended relationship requires review.
\item The registry labels the statement \texttt{corrected\_verified} and reports that the official Oberwolfach report was checked directly and that definitions were restored while extraction artifacts were removed. This is supplied verification metadata, not a claim of source consultation or mathematical verification in this draft. The supplied original excerpt only conjectures that an unspecified map \(\rho\) is a homeomorphism for every finite tensor category.
\end{itemize}
\end{document}
